L'isòtop del sodi ${}^{24}_{11}\mathrm{Na}$ té un excés de neutrons i sabem que és un emissor beta.

\begin{apartat}{1,25}
Raoneu si el sodi 24 serà una font d'emissió d'electrons o de positrons. Escriviu la reacció de la desintegració beta del ${}^{24}_{11}\mathrm{Na}$ sense oblidar els neutrins o antineutrins. Quin és l'altre nucli resultant?
\end{apartat}

\begin{apartat}{1,25}
L'isòtop del sodi 24 té un període de semidesintegració de 15,0~h. Una solució salina amb àtoms de sodi 24 té una activitat inicial de 580~kBq. Aquesta solució s'injecta a un pacient i es dissol per tota la seva sang. Trobeu l'activitat de la solució injectada al cap de 10,0~h. Si al cap d'aquest temps en 1,00~mL de sang hi ha una activitat de 60,0~Bq, quin és el volum de sang del pacient?
\end{apartat}

\dades{%
  \begin{tabular}{|l|c|c|c|c|c|}
    \hline
    \textit{Element} & F & Ne & Na & Mg & Al\\
    \hline
    $Z$ & 9 & 10 & 11 & 12 & 13\\
    \hline
  \end{tabular}}
